% \textbf{Predstavil:}
\textbf{Predstavil:} Gašper Simonič

\begin{enumerate}

\item Naj bo $n$ naravno število.
Z integracijo `po delih' (per partes) pokažite, da velja
    
\[
\int x^n e^{ax} \,dx = \frac{x^n e^{ax}}{a} - \frac{n}{a} \int x^{n-1} e^{ax} \,dx, \quad a \neq 0.
\]
    
\item Pokažite, da velja
    \[
    \int (\ln x)^n \,dx = x (\ln x)^n - n \int (\ln x)^{n-1} \,dx. 
		\]
    
\item S pomočjo zgornjih formul izračunajte naslednje integrale:
		% \begin{multicols}{2}
			\begin{itemize}
				\item $\int x^3 e^{2x} \,dx.$ 
				\item $\int x^2 e^{-x} \,dx.$ 
				\item $\int (\ln x)^3 \,dx.$ 
				\item $\int (\ln x)^4 \,dx.$ 
			\end{itemize}
		% \end{multicols}	
    
\item Kot lahko vidite, če bi integrirali  $\int x^3 e^x \,dx$  z integracijo per partes, bi bil rezultat oblike
		\[
    \int x^3 e^x \,dx = A x^3 e^x + B x^2 e^x + C x e^x + D e^x + E.  
		\]
    
   Odvajajte obe strani te enačbe in izračunajte koeficiente  $A, B, C, D$. Na ta način lahko izračunate integral, ne da bi dejansko izvedli integracijo.
    
\item Dokažite, da če je  $P(x)$  polinom stopnje  $k$ , potem velja
		\[
    \int P(x) e^x \,dx = [P(x) - P'(x) + \dots \pm P^{(k)}(x)] e^x + C.
		\]
    
    \textbf{Namig}: Uporabite lahko podobno idejo kot pri (d)
    
\item Izračunajte naslednje integrale
	\begin{itemize}    
     \item $\int (x^2 - 3x + 1)e^x \,dx.$ 
     \item $\int (x^3 - 2x)e^x \,dx.$
  \end{itemize}
\end{enumerate}